\documentclass[10pt,a4paper]{amsart}
\usepackage{amsmath,amssymb,amsthm,graphicx}
\usepackage{verbatim}
\newenvironment{reference}{\comment}{\endcomment}
%\newenvironment{reference}{}{}
\newenvironment{slogan}{\comment}{\endcomment}
\newenvironment{history}{\comment}{\endcomment}
% For commutative diagrams we use Xy-pic
\usepackage[all]{xy}

% We use 2cell for 2-commutative diagrams.
\xyoption{2cell}
% Theorem environments.
%
\theoremstyle{plain}
\newtheorem{theorem}[subsection]{Theorem}
\newtheorem{proposition}[subsection]{Proposition}
\newtheorem{lemma}[subsection]{Lemma}

\theoremstyle{definition}
\newtheorem{definition}[subsection]{Definition}
\newtheorem{example}[subsection]{Example}
\newtheorem{exercise}[subsection]{Exercise}
\newtheorem{situation}[subsection]{Situation}

\theoremstyle{remark}
\newtheorem{remark}[subsection]{Remark}
\newtheorem{remarks}[subsection]{Remarks}

\numberwithin{equation}{subsection}

% Macros
%
\def\lim{\mathop{\mathrm{lim}}\nolimits}
\def\colim{\mathop{\mathrm{colim}}\nolimits}
\def\Spec{\mathop{\mathrm{Spec}}}
\def\Hom{\mathop{\mathrm{Hom}}\nolimits}
\def\Ext{\mathop{\mathrm{Ext}}\nolimits}
\def\SheafHom{\mathop{\mathcal{H}\!\mathit{om}}\nolimits}
\def\SheafExt{\mathop{\mathcal{E}\!\mathit{xt}}\nolimits}
\def\Sch{\mathit{Sch}}
\def\Mor{\mathop{\mathrm{Mor}}\nolimits}
\def\Ob{\mathop{\mathrm{Ob}}\nolimits}
\def\Sh{\mathop{\mathit{Sh}}\nolimits}
\def\NL{\mathop{N\!L}\nolimits}
\def\CH{\mathop{\mathrm{CH}}\nolimits}
\def\proetale{{pro\text{-}\acute{e}tale}}
\def\etale{{\acute{e}tale}}
\def\QCoh{\mathit{QCoh}}
\def\Ker{\mathop{\mathrm{Ker}}}
\def\Im{\mathop{\mathrm{Im}}}
\def\Coker{\mathop{\mathrm{Coker}}}
\def\Coim{\mathop{\mathrm{Coim}}}

% Boxtimes
%
\DeclareMathSymbol{\boxtimes}{\mathbin}{AMSa}{"02}

%
% Macros for moduli stacks/spaces
%
\def\QCohstack{\mathcal{QC}\!\mathit{oh}}
\def\Cohstack{\mathcal{C}\!\mathit{oh}}
\def\Spacesstack{\mathcal{S}\!\mathit{paces}}
\def\Quotfunctor{\mathrm{Quot}}
\def\Hilbfunctor{\mathrm{Hilb}}
\def\Curvesstack{\mathcal{C}\!\mathit{urves}}
\def\Polarizedstack{\mathcal{P}\!\mathit{olarized}}
\def\Complexesstack{\mathcal{C}\!\mathit{omplexes}}
% \Pic is the operator that assigns to X its picard group, usage \Pic(X)
% \Picardstack_{X/B} denotes the Picard stack of X over B
% \Picardfunctor_{X/B} denotes the Picard functor of X over B
\def\Pic{\mathop{\mathrm{Pic}}\nolimits}
\def\Picardstack{\mathcal{P}\!\mathit{ic}}
\def\Picardfunctor{\mathrm{Pic}}
\def\Deformationcategory{\mathcal{D}\!\mathit{ef}}
\usepackage{cleveref}
\crefname{originallemma}{original lemma}{original lemmas}
\Crefname{originallemma}{Original lemma}{Original lemmas}
\makeatletter
\newlabel{lemma-base-change-gerbe}{{lemma-base-change-gerbe}{source}{}{sourcepointer-1}{}}
\newlabel{lemma-base-change-gerbe@cref}{{[originallemma][0][]lemma-base-change-gerbe}{[1][0][]source}{}{}{}}
\makeatother
\crefname{originallemma}{original lemma}{original lemmas}
\Crefname{originallemma}{Original lemma}{Original lemmas}
\makeatletter
\newlabel{lemma-gerbe-descent}{{lemma-gerbe-descent}{source}{}{sourcepointer-2}{}}
\newlabel{lemma-gerbe-descent@cref}{{[originallemma][0][]lemma-gerbe-descent}{[1][0][]source}{}{}{}}
\makeatother
\crefname{originallemma}{original lemma}{original lemmas}
\Crefname{originallemma}{Original lemma}{Original lemmas}
\makeatletter
\newlabel{lemma-gerbe-with-section}{{lemma-gerbe-with-section}{source}{}{sourcepointer-3}{}}
\newlabel{lemma-gerbe-with-section@cref}{{[originallemma][0][]lemma-gerbe-with-section}{[1][0][]source}{}{}{}}
\makeatother
\crefname{originalreference}{original reference}{original references}
\Crefname{originalreference}{Original reference}{Original references}
\makeatletter
\newlabel{spaces-groupoids-lemma-group-quotient-gerbe}{{spaces-groupoids-lemma-group-quotient-gerbe}{source}{}{sourcepointer-4}{}}
\newlabel{spaces-groupoids-lemma-group-quotient-gerbe@cref}{{[originalreference][0][]spaces-groupoids-lemma-group-quotient-gerbe}{[1][0][]source}{}{}{}}
\makeatother
\crefname{originalreference}{original reference}{original references}
\Crefname{originalreference}{Original reference}{Original references}
\makeatletter
\newlabel{stacks-lemma-when-gerbe}{{stacks-lemma-when-gerbe}{source}{}{sourcepointer-5}{}}
\newlabel{stacks-lemma-when-gerbe@cref}{{[originalreference][0][]stacks-lemma-when-gerbe}{[1][0][]source}{}{}{}}
\makeatother
\title{Locally a gerbe over an algebraic space is a quotient stack}
\author{The Stacks Project Authors}
\date{stacks-project:stacks-morphisms.tex:local-snapshot-20261012; GNU FDL 1.2 or later}
\begin{document}
\maketitle
\noindent\textit{Editorial note.} All delivered mathematical text is the author's own, copied from the pinned Stacks Project chapter source sources/stacks-algebraic-stacks-morphisms/pinned-original-source.tex; no statement or proof step was rewritten, reordered or removed.\par\medskip
\noindent\textit{Editorial note.} Preamble provenance: the Stacks Project keeps one single shared preamble for every chapter; this card takes it directly from the staged sources/stacks-algebraic-stacks-morphisms/preamble.tex as three verbatim spans (lines 8-12, 14-18 and 154-217), so no derived file was written and nothing was added or redefined. The wrapper class and the base packages are supplied by the assembler.\par\medskip
\noindent\textit{Editorial note.} Bibliography provenance: neither the counted statement nor its proof contains a citation (this lane rescans the whole span before writing the card), so the card carries no bibliography and no entry was reconstructed.\par\medskip
\noindent\textit{Editorial note.} Selection provenance: this card was opened under the file-based control contract (CONTROL.md), with the mechanical pre-card gate run on its statement and proof interval before the card was written.\par\medskip
\noindent\textit{Allowed standard prerequisites (editorial).}\begin{itemize}
\item The description of a gerbe with a section as the quotient by the automorphism group of that section. Invoked at the start of the counted proof (source lines 6262-6268). Gives the quotient presentation after producing a section.
\item The existence of a flat finitely presented cover of the base over which the gerbe acquires a section. Invoked inside the counted proof (source lines 6269-6278). Produces the base change on which the presentation is available.
\end{itemize}\medskip
\ifdefined\hypertarget\hypertarget{sourcepointer-1}{}\fi
\noindent\textit{Source reference.} Original lemma lemma-base-change-gerbe: sources/stacks-algebraic-stacks-morphisms/pinned-original-source.tex, source lines 6129--6129. Navigation pointer only: the counted text refers to the Stacks item lemma-base-change-gerbe, which is an existing labeled result of this same chapter and is not part of this card. The pointer carries the target's own label as its number and its definition line as its span; its content is not reproduced here.\par\medskip
\ifdefined\hypertarget\hypertarget{sourcepointer-2}{}\fi
\noindent\textit{Source reference.} Original lemma lemma-gerbe-descent: sources/stacks-algebraic-stacks-morphisms/pinned-original-source.tex, source lines 6161--6161. Navigation pointer only: the counted text refers to the Stacks item lemma-gerbe-descent, which is an existing labeled result of this same chapter and is not part of this card. The pointer carries the target's own label as its number and its definition line as its span; its content is not reproduced here.\par\medskip
\ifdefined\hypertarget\hypertarget{sourcepointer-3}{}\fi
\noindent\textit{Source reference.} Original lemma lemma-gerbe-with-section: sources/stacks-algebraic-stacks-morphisms/pinned-original-source.tex, source lines 6183--6183. Navigation pointer only: the counted text refers to the Stacks item lemma-gerbe-with-section, which is an existing labeled result of this same chapter and is not part of this card. The pointer carries the target's own label as its number and its definition line as its span; its content is not reproduced here.\par\medskip
\ifdefined\hypertarget\hypertarget{sourcepointer-4}{}\fi
\noindent\textit{Source reference.} Original reference spaces-groupoids-lemma-group-quotient-gerbe: sources/stacks-algebraic-stacks-morphisms/pinned-original-source.tex, source lines 6266--6266. Navigation pointer only: the counted text refers to the Stacks item spaces-groupoids-lemma-group-quotient-gerbe, which belongs to another chapter of the project and is not part of this card. The pointer carries the target's own label as its number because this chapter's staged snapshot does not contain the other chapter's tag table; its content is not reproduced here.\par\medskip
\ifdefined\hypertarget\hypertarget{sourcepointer-5}{}\fi
\noindent\textit{Source reference.} Original reference stacks-lemma-when-gerbe: sources/stacks-algebraic-stacks-morphisms/pinned-original-source.tex, source lines 6111--6111. Navigation pointer only: the counted text refers to the Stacks item stacks-lemma-when-gerbe, which belongs to another chapter of the project and is not part of this card. The pointer carries the target's own label as its number because this chapter's staged snapshot does not contain the other chapter's tag table; its content is not reproduced here.\par\medskip
\begin{lemma}
\label{lemma-local-structure-gerbe}
Let $\pi : \mathcal{X} \to \mathcal{Y}$ be a morphism of algebraic stacks.
The following are equivalent
\begin{enumerate}
\item $\mathcal{X}$ is a gerbe over $\mathcal{Y}$, and
\item there exists an algebraic space $U$, a group algebraic space $G$
flat and locally of finite presentation over $U$, and a
surjective, flat, and locally finitely presented
morphism $U \to \mathcal{Y}$ such that
$\mathcal{X} \times_\mathcal{Y} U \cong [U/G]$ over $U$.
\end{enumerate}
\end{lemma}

\begin{proof}
Assume (2). By
Lemma \ref{lemma-gerbe-descent}
to prove (1) it suffices to show that $[U/G]$ is a gerbe over $U$.
This is immediate from
Groupoids in Spaces, Lemma \ref{spaces-groupoids-lemma-group-quotient-gerbe}.

\medskip\noindent
Assume (1). Any base change of $\pi$ is a gerbe, see
Lemma \ref{lemma-base-change-gerbe}.
As a first step we choose a scheme $V$ and a surjective smooth morphism
$V \to \mathcal{Y}$. Thus we may assume that $\pi : \mathcal{X} \to V$
is a gerbe over a scheme. This means that there exists an
fppf covering $\{V_i \to V\}$ such that the fibre category
$\mathcal{X}_{V_i}$ is nonempty, see
Stacks, Lemma \ref{stacks-lemma-when-gerbe} (2)(a).
Note that $U = \coprod V_i \to V$ is surjective, flat, and
locally of finite presentation. Hence we may replace $V$ by $U$ and
assume that $\pi : \mathcal{X} \to U$ is a gerbe over a scheme $U$ and
that there exists an object $x$ of $\mathcal{X}$ over $U$. By
Lemma \ref{lemma-gerbe-with-section}
we see that $\mathcal{X} = [U/G]$ over $U$ for some flat
and locally finitely presented group algebraic space $G$ over $U$.
\end{proof}

\end{document}
